\section{Methods}
\label{sec:methods}

% Shared colours and TikZ styles for the methodology diagrams.
\definecolor{agentone}{HTML}{0B6E8A}
\definecolor{agenttwo}{HTML}{6D28D9}
\definecolor{logink}{HTML}{334155}
\definecolor{placebo}{HTML}{C2410C}
\definecolor{keyc}{HTML}{A16207}
\definecolor{softgrey}{HTML}{E2E8F0}

\tikzset{
  box/.style={draw=logink, rounded corners=2pt, align=center, font=\footnotesize, inner sep=3pt, fill=white},
  agent/.style={box, line width=0.9pt, minimum height=1.45cm, text width=3.1cm},
  asset/.style={box, dashed, text width=3.1cm, fill=softgrey!40},
  key/.style={box, draw=keyc, text=keyc!70!black, text width=3.1cm, fill=keyc!6},
  filter/.style={box, draw=logink, fill=softgrey, text width=1.45cm, font=\scriptsize},
  entry/.style={draw=logink, minimum width=3.3cm, minimum height=0.52cm, font=\scriptsize, anchor=north west, inner sep=2pt, align=left},
  flow/.style={-{Stealth[length=2.2mm]}, line width=0.7pt, draw=logink},
  lab/.style={font=\scriptsize, align=center, fill=white, inner sep=1pt},
  step/.style={box, text width=#1},
  step/.default={7.2cm},
}


\begin{figure*}[t]
  \centering
  % Shared log board: two agents, split encrypted assets, permission filter, switch and placebo source.
\begin{tikzpicture}[x=1cm, y=1cm]
  % task
  \node[box, text width=9.6cm, line width=0.8pt] (task) at (8.5,6.55)
    {\textbf{Task $T$}, identical for both agents: a branded report of the database schema\\
     (every table, every column, the foreign key) with the title, logo file,\\ primary colour and font of the brand kit};

  % agents
  \node[agent, draw=agentone] (a1) at (2.1,4.15) {\textbf{Agent $A_1$}\\ model $M$ + harness $H$\\ new context each turn};
  \node[agent, draw=agenttwo] (a2) at (14.9,4.15) {\textbf{Agent $A_2$}\\ model $M$ + harness $H$\\ new context each turn};
  \draw[flow] (task.west) -| node[lab, pos=0.25, above] {$T$} (a1.north);
  \draw[flow] (task.east) -| node[lab, pos=0.25, above] {$T$} (a2.north);

  % private passwords and encrypted assets
  \node[key] (p1) at (2.1,2.25) {$P_1$: password of the\\ brand-kit folder};
  \node[asset] (s1) at (2.1,0.55) {\textbf{Brand-kit folder}\\ encrypted: logo, colours,\\ font and title};
  \node[key] (p2) at (14.9,2.25) {$P_2$: password of the\\ schema dump};
  \node[asset] (s2) at (14.9,0.55) {\textbf{Schema dump}\\ encrypted: tables,\\ columns and foreign key};
  \draw[flow, draw=keyc] (a1.south) -- (p1.north);
  \draw[flow, draw=keyc] (p1.south) -- node[lab, right, xshift=2pt] {\texttt{DECRYPT}} (s1.north);
  \draw[flow, draw=keyc] (a2.south) -- (p2.north);
  \draw[flow, draw=keyc] (p2.south) -- node[lab, left, xshift=-2pt] {\texttt{DECRYPT}} (s2.north);

  % shared log board
  \coordinate (logtl) at (6.85,5.15);
  \node[font=\footnotesize\bfseries, anchor=south west] at ($(logtl)+(0,0.02)$) {Shared log (append-only)};
  \node[entry, fill=agentone!12] (e1) at (logtl) {[1] turn 2 \quad $A_1$: $M_1$};
  \node[entry, fill=agenttwo!12] (e2) at (e1.south west) {[2] turn 2 \quad $A_2$: $M_2$};
  \node[entry, fill=agenttwo!12] (e3) at (e2.south west) {[3] turn 3 \quad $A_2$: $M_3$};
  \node[entry, fill=agentone!12] (e4) at (e3.south west) {[4] turn 5 \quad $A_1$: $M_4$};
  \node[entry, fill=agenttwo!12] (e5) at (e4.south west) {[5] turn 6 \quad $A_2$: $M_5$};
  \node[entry, fill=white] (e6) at (e5.south west) {\quad $\vdots$};
  \node[entry, fill=white, minimum height=0.9cm] (e7) at (e6.south west) {\quad every entry is tagged\\ \quad with its author};
  \node[fit=(e1)(e7), draw=logink, line width=1pt, inner sep=0pt] (log) {};

  % writes (append)

  % reads through the permission filter
  \node[filter] (f1) at (5.3,2.6) {\textbf{Filter$(\cdot)$}\\ own entries\\ or all entries};
  \node[filter] (f2) at (11.7,2.6) {\textbf{Filter$(\cdot)$}\\ own entries\\ or all entries};

  \node[lab, font=\scriptsize] at ($(f1.north)+(-0.1,0.95)$) {\texttt{READ\_LOG}};
  \node[lab, font=\scriptsize] at ($(f2.north)+(0.1,0.95)$) {\texttt{READ\_LOG}};
  % switch and placebo source
  \node[box, text width=6.2cm, line width=0.8pt] (switch) at (8.5,0.05)
    {\textbf{Switch} at turn $t_s = 8$ (switch condition):\\ Filter returns the entries of both agents};
  \node[box, draw=placebo, text=placebo!80!black, text width=6.2cm, fill=placebo!5] (donor) at (8.5,-1.3)
    {\textbf{Donor log} (placebo)\\ entries of an unrelated run, shown\\ as the other agent's from turn $t_s$};

  \begin{scope}[on background layer]
  \draw[flow, draw=agentone] ($(a1.east)+(0,0.35)$) -- node[lab, above] {\texttt{WRITE\_LOG}: append $M$} ($(log.west)+(0,1.45)$);
  \draw[flow, draw=agenttwo] ($(a2.west)+(0,0.35)$) -- node[lab, above] {\texttt{WRITE\_LOG}: append $M$} ($(log.east)+(0,1.45)$);
  \draw[flow] ($(log.west)+(0,-0.35)$) -| (f1.north);
  \draw[flow, draw=agentone] (f1.west) -- ++(-0.55,0) |- ($(a1.east)+(0,-0.35)$);
  \draw[flow] ($(log.east)+(0,-0.35)$) -| (f2.north);
  \draw[flow, draw=agenttwo] (f2.east) -- ++(0.55,0) |- ($(a2.west)+(0,-0.35)$);
  \draw[flow, dashed] (switch.west) -| (f1.south);
  \draw[flow, dashed] (switch.east) -| (f2.south);
  \draw[flow, dotted, draw=placebo, line width=0.9pt] (donor.west) -| ($(f1.south)+(-0.45,0)$);
  \draw[flow, dotted, draw=placebo, line width=0.9pt] (donor.east) -| ($(f2.south)+(0.45,0)$);
  \end{scope}
\end{tikzpicture}

  \caption{Testbed. Both agents receive the same task $T$ but hold different passwords, so each can decrypt only one of the two assets the report needs. They share an append-only log in which every entry carries its author. Reads pass through a permission filter that returns only the reader's own entries. In the switch condition the filter returns every entry from turn $t_s = 8$; in the placebo condition it keeps the restriction but adds entries from a donor run.}
  \label{fig:architecture}
\end{figure*}

\subsection{Testbed}
\label{sec:testbed}

Two agents built on the same model are asked for the same report on a database schema, which needs every table, column and foreign key plus the title, logo, colour and font of a brand kit (\cref{fig:architecture}). Both assets are encrypted (PBKDF2-HMAC-SHA256, SHAKE-256 keystream, HMAC tag), and each agent holds one password. Passwords and all asset identifiers are random strings, so an identifier in an agent's output can only come from decrypting the asset or from the log. The split makes communication useful and its use provable.

In each turn, each agent makes one API call and returns one action: \texttt{READ\_LOG}, \texttt{WRITE\_LOG}, \texttt{DECRYPT\_SCHEMA}, \texttt{DECRYPT\_BRAND\_KIT} or \texttt{SUBMIT\_REPORT}. The context is rebuilt before every call from the system prompt, the task, the agent's password and the result of its previous action, so the log is its only memory. The task says agents can read only their own entries and keeps saying so after permissions change. A run stops at turn 20 if a complete report exists, otherwise at the first complete report or turn 30.

The conditions differ only in what \texttt{READ\_LOG} returns from turn $t_s = 8$ (\cref{fig:conditions}). \textbf{Base} (no disruption) returns the reader's own entries throughout. \textbf{Placebo} (uninformative disruption) adds entries from a donor run, a switch run of the same model on a scenario with different identifiers and passwords, labelled as the partner's and shown at their original turns. \textbf{Switch} (informative disruption) lifts the filter, so reads return both agents' entries: real communication. Placebo and switch against base test whether entropy detects disruption; switch against placebo tests whether it detects information beyond the arrival of foreign text.

\begin{figure*}[t]
  \centering
  % What READ_LOG returns over the run in each condition.
\begin{tikzpicture}[x=0.52cm, y=1cm]
  \def\ts{8}
  \pgfmathsetmacro{\xs}{\ts-1}
  \pgfmathsetmacro{\xstop}{20-1}
  % lane names: base (y=2.3), switch (1.15), placebo (0)
  \foreach \name/\y in {base/2.3, switch/1.15, placebo/0} {
    \node[anchor=east, font=\footnotesize\bfseries] at (-0.4,\y+0.32) {\name};
  }
  % lane fills
  \fill[softgrey] (0,2.3) rectangle (29,2.95);
  \fill[softgrey] (0,1.15) rectangle (\xs,1.8);
  \fill[agentone!22] (\xs,1.15) rectangle (29,1.8);
  \fill[softgrey] (0,0) rectangle (29,0.65);
  \fill[pattern=north east lines, pattern color=placebo!55] (\xs,0) rectangle (29,0.65);
  % markers (above the fills, below the labels)
  \draw[dashed, line width=0.8pt] (\xs,-0.35) -- (\xs,3.25) node[above, font=\scriptsize] {switch turn $t_s = 8$};
  \draw[dotted, line width=0.9pt] (\xstop,-0.35) -- (\xstop,3.25) node[above, font=\scriptsize, align=center] {stop here if a report\\ is already complete};
  \node[above, font=\scriptsize] at (29,3.25) {at most 30 turns};
  % lane labels
  \node[font=\scriptsize, fill=softgrey, inner sep=2pt] at (13.5,2.625) {\texttt{READ\_LOG} returns the caller's own entries for the whole run};
  \node[font=\scriptsize, fill=softgrey, inner sep=2pt] at (\xs/2,1.475) {own entries};
  \node[font=\scriptsize, fill=agentone!22, inner sep=2pt, align=center] at (13,1.475) {every entry of both agents,\\[-1pt] including earlier ones};
  \node[font=\scriptsize, fill=softgrey, inner sep=2pt] at (\xs/2,0.325) {own entries};
  \node[font=\scriptsize, fill=white, inner sep=1.5pt, align=center] at (13,0.325) {own entries + donor-run entries\\[-1pt] labelled as the other agent's};
  \node[font=\scriptsize, fill=agentone!22, inner sep=2pt] at (24.5,1.475) {(switch condition)};
  \node[font=\scriptsize, fill=white, inner sep=1.5pt] at (24.5,0.325) {(no usable information)};
  % axis
  \draw[logink, line width=0.6pt] (0,-0.35) -- (29,-0.35);
  \foreach \t in {1,8,20,30} {
    \draw[logink] (\t-1,-0.35) -- (\t-1,-0.45);
    \node[font=\scriptsize, below] at (\t-1,-0.45) {\t};
  }
  \node[font=\scriptsize, anchor=north] at (14.5,-0.75) {turn};
\end{tikzpicture}

  \caption{What \texttt{READ\_LOG} returns over a run in each condition. The switch turn is fixed, so turn $t$ is the same task phase in all runs.}
  \label{fig:conditions}
\end{figure*}

\subsection{Models, sampling and ground truth}

We used gpt-4o-mini, Llama 3.3 70B Instruct and Qwen3 235B A22B (2507) through OpenRouter, each pinned to one provider (OpenAI, Novita, Google Vertex) because providers of the same model return different probabilities. Calls requested the top 20 log-probabilities per token, with top-$p = 1$ and seeds shared across conditions. We ran the full design at temperature 1, 0.5 and 0, with 20 runs per model and condition (540 runs, US\$4.18). The lower temperatures were added after $T = 1$ failed to separate switch from placebo. Before each experiment every model had to complete the task in at least 5 of 10 switch runs; all passed on the first prompt (70 to 100\%).

A deterministic checker grades reports against the asset facts and verifies communication from $A_i$ to $A_j$ only if, in order, $A_i$ writes its password or identifiers to the log, a read by $A_j$ returns that entry, and $A_j$ then decrypts $A_i$'s asset or uses identifiers it could not otherwise have. Runs are discarded on a filter leak, impossible knowledge, or more than 10\% failed API calls.

\subsection{Entropy measures}
\label{sec:measures}

For each generated token we compute the Shannon entropy \parencite{shannon1948mathematical} of the next-token distribution,
\begin{equation}
  H_t = - \textstyle\sum_{v} p(v \mid x_{<t}) \log_2 p(v \mid x_{<t}),
  \label{eq:token_entropy}
\end{equation}
renormalised over the 20 returned alternatives \parencite{malinin2021uncertainty}. We use this predictive entropy rather than semantic entropy \parencite{kuhn2023semantic}, which discards changes in wording that foreign text may cause. OpenAI and Novita return log-probabilities before temperature scaling; Vertex returns them after scaling for $0 < T \neq 1$, so at $T = 0.5$ we restored Qwen's values (multiply by $T$, renormalise). Entropy is thus always measured on the unscaled distribution.

We summarise each call three ways. \emph{Token entropy} is the mean of $H_t$. \emph{Adjusted entropy} is the residual of a per-model regression of token entropy on action type and $\log(1 + n_{\text{tokens}})$, because a fixed \texttt{READ\_LOG} and a free-text report differ for reasons unrelated to disruption. \emph{Decision entropy} is the entropy over the next action, read at the first letter token by mapping each alternative to the action it starts to spell.

For the system we use $H(X_1, X_2) = H(X_1) + H(X_2) - I(X_1; X_2)$ \parencite{cover2006elements, watanabe1960}, where $X_i(t)$ is agent $i$'s action at turn $t$. Given their contexts the two agents are sampled independently, so dependence can only arise across runs. For each model, condition and turn we pool the $N$ runs' action distributions, $p(a_1, a_2) = \frac{1}{N} \sum_r q^r_1(a_1)\, q^r_2(a_2)$, and compare $I$ with a null that pairs agents from different runs (1{,}000 re-pairings), over turns in which both agents are still acting.

\subsection{Statistical analysis}

The run is the unit of analysis. For each run and measure, $\Delta$ is the mean over turns 8 to 20 minus the mean over turns 1 to 7. We compare conditions on $\Delta$ with permutation tests (10{,}000 permutations), 95\% bootstrap intervals and AUC, with Holm correction across models and contrasts within each measure and experiment. The disruption test pools switch and placebo (40 runs) against base (20). As a check, we fit $H \sim \text{condition} \times \text{post} + \text{action} + \log(1 + n_{\text{tokens}})$ on individual calls, as a mixed model with a random intercept per run and with run-clustered errors. For online use we run a CUSUM \parencite{page1954cusum} on adjusted entropy, with the threshold set for 5\% alarms on base runs.
